\documentclass[12pt]{article}

\usepackage[a4paper,margin=1in]{geometry}
\usepackage{amsmath, amssymb}
\usepackage{setspace}

\setstretch{1.15}

\title{CSE400 -- Lecture 7 Lecture Scribe\\
Expectation, CDFs, PDFs and Problem Solving}
\date{}

\begin{document}

\maketitle

\section{Lecture Coverage (From Outline Slide)}

The lecture includes:

\begin{enumerate}
\item Cumulative Distribution Function (CDF)
\begin{itemize}
\item Definition
\item Properties
\item Example
\end{itemize}

\item Probability Density Function (PDF)
\begin{itemize}
\item Definition
\item Properties
\item Example
\end{itemize}

\item Expectation of Random Variables
\begin{itemize}
\item Definition and example
\item Expectation of a function of RV
\item Linear operation with expectation
\item n-th moments and central moments (Variance, Skewness, Kurtosis)
\end{itemize}
\end{enumerate}

\section{Cumulative Distribution Function (CDF)}

\subsection{Definition and Notation}

Let $X$ be a random variable.

The cumulative distribution function (CDF) is defined as:
\[
F_X(x) = \Pr(X \le x), \qquad -\infty < x < \infty
\]

Notation

$X$ : random variable

$F_X(x)$ : cumulative distribution function of $X$

Statement from slides:

Most information about the random experiment described by $X$ is determined by the behavior of $F_X(x)$.

\subsection{Assumptions and Conditions (Properties of a CDF)}

A function must satisfy the following conditions to be a valid CDF.

\begin{enumerate}
\item Bounds
\[
0 \le F_X(x) \le 1
\]

\item Limits
\[
F_X(-\infty) = 0, \qquad F_X(\infty) = 1
\]

\item Monotonicity

If $x_1 < x_2$, then
\[
F_X(x_1) \le F_X(x_2)
\]

\item Probability from CDF

For $x_1 < x_2$,
\[
\Pr(x_1 < X \le x_2) = F_X(x_2) - F_X(x_1)
\]
\end{enumerate}

\section{Worked Example 1 --- Validity of CDF}

Problem: Determine which functions are valid CDFs.

\subsection*{(a)}

\[
F_X(x) = \frac{1}{2} + \frac{1}{\pi}\tan^{-1}(x)
\]

Step-by-step verification

Step 1: Check limits.

As $x \to -\infty$,
\[
\tan^{-1}(x) \to -\frac{\pi}{2}
\]
\[
F_X(-\infty) = \frac{1}{2} - \frac{1}{2} = 0
\]

As $x \to \infty$,
\[
\tan^{-1}(x) \to \frac{\pi}{2}
\]
\[
F_X(\infty) = \frac{1}{2} + \frac{1}{2} = 1
\]

Step 2: Since $\tan^{-1}(x)$ is increasing, the function is non-decreasing.

Step 3: Values remain between 0 and 1.

✔ Valid CDF.

\subsection*{(b)}

\[
F_X(x) = (1 - e^{-x})u(x)
\]

Step-by-step verification

Step 1: For $x < 0$,
\[
u(x) = 0 \Rightarrow F_X(x) = 0
\]

Step 2: For $x \ge 0$,
\[
F_X(x) = 1 - e^{-x}
\]

Step 3: Limits:
\[
F_X(-\infty) = 0, \qquad F_X(\infty) = 1
\]

Step 4: Function increases with $x$.

✔ Valid CDF.

\subsection*{(c)}

\[
F_X(x) = e^{-x^2}
\]

As $x \to \infty$,
\[
F_X(x) \to 0
\]
which violates $F_X(\infty) = 1$.

✘ Invalid CDF.

\subsection*{(d)}

\[
F_X(x) = x^2 u(x)
\]

As $x \to \infty$,
\[
F_X(x) \to \infty
\]
which violates $0 \le F_X(x) \le 1$.

✘ Invalid CDF.

\section{Worked Example 2 --- Computing Probabilities from a CDF}

Given:
\[
F_X(x) = (1 - e^{-x})u(x)
\]

\subsection*{(a)}

\[
\Pr(X > 5)
\]

Step 1:
\[
\Pr(X > 5) = 1 - \Pr(X \le 5)
\]

Step 2:
\[
= 1 - F_X(5)
\]

Step 3:
\[
= 1 - (1 - e^{-5})
\]

Result:
\[
= e^{-5}
\]

\subsection*{(b)}

\[
\Pr(X < 5)
\]

\[
\Pr(X < 5) = F_X(5) = 1 - e^{-5}
\]

\subsection*{(c)}

\[
\Pr(3 < X < 7)
\]

Step 1:
\[
\Pr(3 < X < 7) = F_X(7) - F_X(3)
\]

Step 2:
\[
= (1 - e^{-7}) - (1 - e^{-3})
\]

Result:
\[
= e^{-3} - e^{-7}
\]

\subsection*{(d) Conditional Probability}

\[
\Pr(X > 5 \mid X < 7)
\]

Step 1:

Let
\[
A = \{X > 5\}, \qquad B = \{X < 7\}
\]

Step 2:
\[
\Pr(A \mid B) = \frac{\Pr(A \cap B)}{\Pr(B)}
\]

\[
\Pr(A \cap B) = \Pr(5 < X < 7) = F_X(7) - F_X(5)
\]

\[
\Pr(B) = F_X(7)
\]

Step 3:

Final expression:
\[
\Pr(X > 5 \mid X < 7)
=
\frac{F_X(7) - F_X(5)}{F_X(7)}
\]

\section{Probability Density Function (PDF)}

\subsection{Definition}

For a continuous random variable $X$,

\[
f_X(x) =
\lim_{\epsilon \to 0}
\frac{\Pr(x \le X < x+\epsilon)}{\epsilon}
\]

\subsection{Relationship Between PDF and CDF (Step-by-Step Derivation)}

From the slides:

For a continuous range,
\[
\Pr(x \le X < x+\epsilon) = F_X(x+\epsilon) - F_X(x)
\]

Substitute into PDF definition:
\[
f_X(x) =
\lim_{\epsilon \to 0}
\frac{F_X(x+\epsilon) - F_X(x)}{\epsilon}
\]

By definition of derivative,
\[
f_X(x) = \frac{d}{dx}F_X(x)
\]

\subsection{Result (Converse Statement)}

The CDF can be written as the integral of the PDF:
\[
F_X(x) = \int_{-\infty}^{x} f_X(t)\, dt
\]

\section{Expectation of Random Variables (As Listed in Slides)}

The lecture includes:

Definition and example of expectation

Expectation of a function of a random variable

Linear operation with expectation

n-th moments and central moments:
\begin{itemize}
\item Variance
\item Skewness
\item Kurtosis
\end{itemize}

(Only topics listed in slides are included; no additional formulas introduced.)

End of Lecture Scribe

This lecture scribe is:

Fully restricted to lecture slide content

Structured for exam revision

Step-by-step where derivations appear in slides

Free from external intuition or added explanations

If required, another independent version with a different structural ordering can be generated while remaining fully slide-faithful.

\end{document}
